\documentclass[11pt,a4paper,oneside]{article}

\usepackage[a4paper,left=3cm,right=2cm,top=2.5cm,bottom=2.5cm]{geometry}
\usepackage{tabto}
\usepackage[french]{babel}
\usepackage[T1]{fontenc}
\usepackage[utf8]{inputenc}
\usepackage{standalone}

\usepackage[
backend=biber,
style=alphabetic,
sorting=ynt
]{biblatex}
\addbibresource{../manuscript/references/references.bib}

% \input{CV-short/structure.tex} % Include the file specifying document layout and packages

\begin{document}

\begin{titlepage}

	\vspace*{5.5cm}

	\begin{center}

		\noindent {\Huge Dossier de candidature à un poste d'Attaché Temporaire d'Enseignement et de Recherche à l'Université de Poitiers}

		\vspace{1.5cm}

		{\huge Section 27}

		\vspace{2.5cm}

		\noindent {\LARGE Maxime Gaide}

		\vspace{1.5cm}

		\noindent {\Large \today}

	\end{center}

	\vfill

	\noindent Doctorant au Laboratoire d'Informatique et d'Automatique pour les Systèmes

	\noindent ISAE-ENSMA / Université de Poitiers

	% \noindent 2024
\end{titlepage}

\newpage

\section{Informations générales}

\noindent Nom Complet : Maxime Gaide

\noindent Âge : 27 ans

\noindent Situation actuelle : Doctorant en troisième année

\noindent Laboratoire : \tabto{2.35cm}Laboratoire d'Informatique et d'Automatique pour les Systèmes (LIAS)

\noindent \tabto{2.35cm}  ISAE-ENSMA / Université de Poitiers

\noindent Thématique de recherche : Informatique graphique analyse et reconstruction de formes

\noindent Section CNU : 27

\noindent Adresse postale : 4 IMP DE LA HAUTE PAYRE, 86130 Jaunay-Marigny, France

\noindent Téléphone : +33786269524

\noindent Courriel : maxime.gaide@ensma.fr

\noindent Langues :
\begin{itemize}
	\item français : langue maternelle
	\item anglais : professionnel scientifique écrit et oral
	\item suédois : débutant
	\item allemand : débutant
\end{itemize}

\section{Formation Initiale}

\subsection{Licence Informatique - Aix-Marseille Université (2014-2018)}

La première année de licence permet d'acquérir les fondamentaux de
l'informatique avec des bases en mathématiques. %
Les modules d'introduction à la programmation sont enseignés en \textbf{langage
	C}. %
Les deux années suivantes forment, entre autres, à l'algorithmique, la théorie
des graphes, la programmation orienté objet (\textbf{Java}) ou encore à la
compréhension des systèmes informatiques. %
Cette formation offre également la possibilité de personnaliser ses compétences
via des modules optionnels (programmation fonctionnelle \textbf{OCaml},
connaissance de la recherche en informatique). %

\subsection{Master 1 Informatique - Université d'Uppsala, Suède, (2018-2019)}

Cette année a été réalisée en Suède via une mobilité Erasmus+ comprenant un test
d'anglais au début ainsi qu'à la fin. %
Pour cette année, j'ai suivi un cours intensif de suédois (niveau A1), des
modules d'introduction à la programmation parallèle (\textbf{pthread},
\textbf{OpenMP} et \textbf{OpenMPI}), d'interaction Homme-Machine, de base de
données, gestion de projet et d'introduction à l'informatique graphique. %

\subsection{Stage recherche - Détection de cratères dans les maillages
	surfaciques de corps célestes - Aix-Marseille Université (M1, 2019, 2 mois)}

Ce stage de recherche a eu lieu entre le Laboratoire d'Informatique et Système
(LIS) d'Aix-Marseille Université et le Laboratoire d'Astrophysique de
Marseille. %
L'objectif était, dans un premier temps, de concevoir et développer un outil
permettant de détecter des cratères~; et dans un second temps, d'extraire des
caractéristiques à partir des cratères détectés. %
Ces caractéristiques comprennent entre autre le diamètre du cratère, sa
profondeur et son type. %
Au cours cette première expérience de recherche, j'ai assisté un autre étudiant
(en M2) pour développer un outil répondant aux problématiques du sujet. %

\subsection{Master 2 Géométrie et Informatique Graphique - Aix-Marseille Université (2019-2020)}

La deuxième année se spécialise dans la modélisation 3D, l'informatique
graphique et la géométrie appliquée. %
Le semestre 3 forme aux fondamentaux de la géométrie discrète, de l'animation, à
l'analyse et au traitement de la géométrie. %
Ce parcours de Master étant indifférencié vis-à-vis d'un parcours recherche ou
professionnel propose également un module autour de la réalisation d'un projet
de type industriel avec un cahier des charges et sa documentation. %
Le semestre 4 comporte un projet de fin d'étude ainsi qu'un stage de 6 mois. %
Mon projet de fin d'étude concernait la squelettisation duale-primale à partir
de maillages simpliciaux 3-dimensionnels (utilisation de \textbf{VTK}). Mon
stage était un stage recherche sur la compression d'information géométrique pour
la diffusion d'un flux de rendu graphique. %

\subsection{Projet de fin d'études - Squelettisation duale-primale à partir de
	maillages simpliciaux 3-dimensionels (M2, 2020, 2 mois)}

Ce projet, orienté recherche, a eu lieu sous la supervision de Jean-Luc Mari et Ricardo
Uribe-Lobello membres de l'équipe Géométrie et Modélisation du LIS. %
L'objet de ce projet était de vérifier si les critères de squelettisation
duale-primale proposés dans \cite{URI2019Dual} étaient suffisamment robustes
pour être généralisées en toute dimension. %
Pour cela, avec mon groupe, nous avons procédé à la réimplantation des méthodes
proposées en utilisant les outils de la suite \textbf{VTK}.

\subsection{Stage recherche - Compression de pipeline 3D pour le streaming d'application - Nothing2Install, Marseille, (M2, 2020, 6 mois)}

Nothing2Install est une Start-Up proposant une solution permettant l'utilisation
d'applications sans la contrainte de les installer. %
Le principe de cette solution est la transmission des instructions de rendu
graphique d'une application depuis un serveur vers un client. %
Cependant, un flux de rendu graphique contient, entre autre de la géométrie, de
la texture, etc. %
Toutes ces données accumulées sont susceptibles de saturer la bande passante et
impacter négativement le réseau. %
Au cours de ce stage, j'ai travaillé sur la compression d'information
géométrique entre deux cycles de rendu. %
Mon objectif était de compresser des informations géométriques et comparer avec
des outils de compression plus généralistes (Zstandard, GZip, …) voire de
composer avec. %
L'objectif étant d'alléger la charge de diffusion sur un réseau. %

\subsection{Doctorat en Informatique - Rejeu et modélisation spatio-temporelle à base de règles d'éléments archéologiques  (2021-en cours)}

\noindent Directeur de thèse~: Stéphane Jean

\noindent Encadrants~: \tabto{2.26cm} David Marcheix, Xavier Skapin, Agnès Arnould et Hakim Belhaouari

\noindent Employeur~: ISAE-ENSMA avec financement Ministère

Cette thèse sous contrat doctorale est liée au domaine de la CAO et plus
particulièrement à la reconstruction automatique de modèles 3D. %
La spécificité de ces travaux est de s'appuyer sur des opérations de
modélisation définies à partir de règles de transformation de graphes Jerboa \cite{belhaouari2014jerboa}. %
L'objectif final étant de permettre la génération rapide de variations à partir
d'un modèle initial. %

Les objectifs atteints sont~:
\begin{itemize}
	\item l'enregistrement du processus constructif (spécification paramétrique)
	      d'un objet initial,
	\item l'implantation d'un mécanisme pour la réévaluation d'une spécification
	      paramétrique,
	\item l'édition d'une spécification paramétrique (ajout, retrait, déplacement
	      d'opération) et l'automatisation de sa réévaluation pour des règles de transformation de graphe,
\end{itemize}
Les objectifs en cours et à venir sont~:
\begin{itemize}
	\item l'adaptation d'une spécification paramétrique et l'automatisation de sa
	      réévaluation prenant en compte les scripts de règles,
	\item l'animation de l'évolution d'un modèle au cours du temps,
	\item l'annotation de modèles à partir de données historiques hétérogènes.
\end{itemize}

\subsection{Récapitulatif des compétences acquises}

Tableau récapitulatif compétences acquises :
\begin{center}
	\begin{tabular}{|l|l|l|}
		\hline
		Langages          & outils            & Compétences                                 \\
		\hline
		C/C++             &                   & Algorithmique, structure de données,        \\
		–                 &                   & gestion mémoire, programmation graphique    \\
		–                 & OpenMP, OpenMPI   & Paraléllisation de tâche,                   \\
		–                 & OpenMesh, OpenGL, & Analyse et traitement de maillage,          \\
		–                 & OpenCV            & visualisation, animation et rendu d'objet,  \\
		–                 & –                 & traitement d'image                          \\
		\hline
		Java              &                   & Conception et gestion de projet             \\
		–                 &                   & tests unitaires                             \\
		–                 & JavaFX            & Interface utilisateur                       \\
		-                 & Jerboa            & Conception et développement de règles       \\
		                  &                   & de transformations de graphes               \\
		                  &                   & modélisation géométrique à base topologique \\
		\hline
		OCaml et Poly/SML &                   & Bases de la programmation fonctionnelle     \\
		\hline
		Python            & Pandas, Numpy     & Prototypage, développement outil, scripts   \\
		–                 & Graphviz          & visualisation de structure                  \\
		\hline
		                  & Bash              & Scripts et outils systèmes                  \\
		\hline
		                  & LaTeX             & Rédaction de documents et présentations     \\
		\hline
		                  & Git               & Contrôle de version de projet               \\
		\hline
	\end{tabular}
\end{center}
\section{Recherche}

\subsection{Analyses syntaxiques de règles de transformations de graphe}
La plupart des logiciels de modélisation géométrique disposent d'opérations
développées sous forme de code. %
Ceci rend leurs analyses particulièrement complexes. %
Ainsi, déterminer quelle est la partie modifiée d'une modèle à partir d'une
opération est difficile voire impossible. %
Deux méthodes sont généralement utilisées pour déterminer l'impacte d'une
opération sur un modèle. %
La première consiste à parcourir entièrement le modèle et à enregistrer les
changements. %
Dans ce cas, la méthode est efficace mais le coût en temps dépend de la taille
du modèle. %
La seconde consiste à ajouter des instructions directement dans l'opération. %
Le coût en temps est négligeable mais le risque et d'introduire des informations
erronées sans pour autant délimiter précisément le périmètre d'application de
l'opération sur le modèle. %
Les règles de transformation de graphes, elles, suivent une méthode formelle. %
L'analyse de leur syntaxe permet de déduire quelle partie d'un objet et impactée
et dans quelle mesure. %
Nous distinguons six types d'évènements topologiques : la création, la
suppression, la modification, le non-changement, la scission ou encore la fusion
de tout type de cellule topologique (sommet, face, arête, etc.). %
Cette analyse est automatique. %
Il est possible, d'une part, de détecter de manière statique l'ensemble des
évènements provoqués par l'application d'une règle, peu import le modèle. %
D'autre part, lorsque la règle est conçue de telle sorte qu'il y a une
incertitude, une vérification localisée permet de confirmer ou non ses
potentielles conséquences. %
Ces travaux ont été publiés dans \cite{gaide2023automatic}. %

\subsection{Réévaluation de modèle à partir de règles de transformation de graphe}
La réévaluation d'un modèle est une fonctionnalité essentielle dans le domaine
de la CAO. %
Cette fonctionnalité est liée à un problème de longue date qui est
l'identification persistante des cellules topologiques dans un modèle
géométrique. %
Afin de procéder à la réévaluation de modèles, nous avons conçu un système de
nommage persistant à partir des règles de transformations Jerboa
\cite{belhaouari2014jerboa}. %
Ce mécanisme de nommage persistant s'appuie un historique des éléments du
modèle, généré à partir des règles, et se propage dans le modèle au fur et à
mesure de sa construction. %
Une fois le modèle construit et sa spécification paramétrique enregistrée,
chacune des entités topologiques étant désigné comme lieu d'application d'une
règle est désigné par un nom persistant. %
Puisqu'un nom persistant est une sorte d'historique, il est alors possible de
décrire l'histoire d'une entité topologique, ses différents états et ses
origines. %
Au cours de la réévaluation, nous pouvons ainsi reconstruire l'objet pas à pas
en mettant à jour l'histoire de chaque entité topologique en fonction de
l'édition de la spécification paramétrique. %
Ces travaux ont été présentés sous forme de poster dans \cite{gaide2023model}. %
Un article est en cours de soumission. %

\subsection{Exploitation des scripts de règles pour la réévaluation de modèles}
Bien que les règles de transformations de graphe s'avèrent êtres des outils
puissants, celles-ci doivent être composées pour définir des opérations
complexes. %
Jerboa \cite{belhaouari2014jerboa} permet de créer des scripts de règle au sein
desquelles il est possible d'appeler des règles, voire d'autres scripts et de
faire usage de structures de contrôle. %
Dans le cadre de ces scripts, nous étudions trois types de structures :
\begin{itemize}
	\item les structures conditionnelles pouvant faire varier un modèle selon son
	      évaluation,
	\item les structures itératives de type «~pour chaque sous-cellule d'une
	      cellule~» dont le nombre d'itérations peut varier et faire varier les
	      noms persistants au rejeu,
	\item les séquences d'applications.
\end{itemize}
À ce jour, ces travaux sont en cours d'étude.

\section{Enseignements et encadrements}

Les enseignements suivants ont été faits dans le cadre d'un Contrat Doctoral avec
Enseignement à l'ISAE-ENSMA et valorisés en équivalent Heure TD.

\subsection{Algorithmique et Structures Numériques; 1ère année ISAE-ENSMA; 27h~TP, 45h~Projet, 3h~TD}
Ce module est enseigné avec le langage \textbf{Ada}. %
L'objectif de cet enseignement est d'acquérir les bases de la programmation. %
Les étudiants apprennent à utiliser les différentes structures de contrôle, à
choisir des structures de données selon leurs besoins et à les manipuler. %
Le langage \textbf{Ada} imposant un typage fort, l'accent est mis sur la
fiabilité des programmes développés. %
Pour ce module, j'ai encadré des séances de TP, Projet (étalé sur 3 séances) et quelques TDs. %
J'ai également pu surveiller des sessions de contrôle continu. %

\subsection{Base de la Conception Logicielle; 1ère année ISAE-ENSMA; 42h~TP}
Dans ce module, les étudiants apprennent à concevoir des programmes respectant
différents concepts tels que la ré-utilisabilité et la maintenabilité du code. %
Ils apprennent également à développer des programmes en utilisant les concepts
de généricité et de polymorphisme.
Pour ce module, j'ai encadré des séances de TP.

\subsection{Utilisation et Exploitation des Données; 1ère année ISAE-ENSMA; 6h~TP}
Ce module est enseigné avec le langage \textbf{Python} et s'appuie sur
l'utilisation de notebooks. %
L'objectif de cet enseignement est de donner les clés pour s'approprier des
connaissances grâce à l'étude de jeux de données. %
Les étudiants apprennent à utiliser différents outils tels que \textbf{pandas},
\textbf{numpy}, \textbf{SQL} pour exploiter des jeux de données et, entre autre,
les visualiser. %
Pour ce module, j'ai pu encadrer des séances de TP.

\subsection{Apprentissage Automatique appliqué; 2ème année ISAE-ENSMA; 3h~TP}
Ce module est enseigné avec le langage \textbf{Python} et s'appuie sur
l'utilisation de notebooks Jupyter. %
L'objectif de cet enseignement est de proposer une introduction à
l'apprentissage automatique. %
Les étudiants acquièrent des connaissances sur les différentes catégories
d'apprentissages, leurs avantages et inconvénients. %
Ils sont également confrontés aux problématiques liés à la validation des
modèles. %
Pour ce module, j'ai encadré une séance de TP. %

\subsection{Encadrement EUR - Conception d'une interface pour le rejeu et
	animation spatio-temporelle de modèle (en cours)}
L'Université de Poitiers propose un parcours EUR pour son master informatique. %
Ce parcours a deux objectifs : premièrement, offrir aux étudiants du master une
expérience à la recherche au sein d'une équipe sur un véritable sujet de
recherche~; deuxièmement, offrir à des doctorants une expérience
d'encadrement. %
Dans ce contexte, j'encadre un étudiant travaillant sur une problématique issue
de mes travaux de thèse. %
À l'issue de cet encadrement, l'étudiant doit rédiger un rapport, concevoir un
poster et présenter son travail. %

\subsection{Récapitulatif des enseignements}

Tableau récapitulatif des enseignements, des niveaux d'études et des heures :
\begin{center}
	\begin{tabular}{|l|l|l|l|l|l|}
		\hline
		Établissement                      & Matières enseignées                     & Niveaux      & \multicolumn{3}{l|}{Nombre d'heures}                \\
		\hline
		                                   &                                         &              & TD                                   & tp  & projet \\
		\hline
		ISAE-ENSMA                         & Algorithmique et Systèmes Numériques    & A1 ingénieur & 3h                                   & 27h & 45h    \\
		\hline
		ISAE-ENSMA                         & Base de la Conception Logicielle        & A1 ingénieur &                                      & 42h &        \\
		\hline
		ISAE-ENSMA                         & Utilisation et Exploitation des Données & A1 ingénieur &                                      & 6h  &        \\
		\hline
		ISAE-ENSMA                         & Apprentissage Automatique Appliqué      & A2 ingénieur &                                      & 3h  &        \\
		\hline
		\multicolumn{3}{|l|}{Total HETD~:} & \multicolumn{3}{l|}{126~h}                                                                                   \\
		\hline
	\end{tabular}
\end{center}

\subsection{Récapitulatif des compétences enseignées}

Tableau récapitulatif compétences enseignées par langage et par niveau :
\begin{center}
	\begin{tabular}{|l|l|l|}
		\hline
		Langages  et outils  & Compétences                                 & Niveaux      \\
		\hline
		Ada                  & Algorithmique, structures de données,       & A1 ingénieur \\
		                     & structures de contrôle, fiabilité du code   &              \\
		                     & conception logicielle, encapsulation        &              \\
		                     & réutilisabilité, généricité, multitâche     &              \\
		\hline
		Python~: Pandas, SQL & Lecture et traitement des données           & A1 ingénieur \\
		- Numpy              & Calcul et statistiques sur les données      &              \\
		- Seaborn            & Visualisation des données                   &              \\
		\hline
		Python~:Scikit-Learn & Apprentissage Automatique, supervisé        & A2 ingénieur \\
		                     & et non supervisé, pré-traitement de donnés, &              \\
		                     & validation de modèle                        &              \\
		\hline
	\end{tabular}
\end{center}

\section{Responsabilités collectives}

\subsection{Étudiant Élu au Conseil du Département Informatique et Interaction, Aix-Marseille Université (L2-L3, 2016-2018) }
J'ai été élu par les autres étudiants d'Aix-Marseille Université pour les
représenter aux conseils du Département Informatique et Interaction.

\subsection{Doctorant Élu au conseil de l'École Doctorale MIMME, Poitiers (2022-2024)}
J'ai été élu par les autres doctorants inscrits auprès de l'ED MIMME pour les
représenter aux conseils de l'ED.

\section{Motivations pour un poste d'ATER à l'Université de Poitiers}

Au cours de mon doctorat, j'ai pu encadrer des étudiants apprenant les bases de
la programmation et de la conception logicielle à l'ISAE-ENSMA. %
Avec cette candidature, je souhaite poursuivre une activité d'enseignement dans
le parcours informatique de l'Université de Poitiers. %
Côté enseignement, cela me permettrait de dispenser de nouveaux enseignements,
de nouvelles compétences et d'en acquérir moi-même. %
Côté recherche, prévoyant de soutenir avant décembre 2024, un poste d'ATER me
permettrait de disposer du temps nécessaire pour finaliser et publier les
travaux en cours. %

\printbibliography

\end{document}
